\documentclass[conference]{IEEEtran}
\IEEEoverridecommandlockouts

\usepackage{cite}
\usepackage{amsmath,amssymb,amsfonts}
\usepackage{algorithmic}
\usepackage{graphicx}
\usepackage{textcomp}
\usepackage{xcolor}
\usepackage{booktabs}
\usepackage{hyperref}
\usepackage{url}

\def\BibTeX{{\rm B\kern-.05em{\sc i\kern-.025em b}\kern-.08em
    T\kern-.1667em\lower.7ex\hbox{E}\kern-.125emX}}

\begin{document}

\title{Honest Ablation of a Hybrid Quantum-Classical Classifier for Early Disease Detection: When Does Training the Quantum Layer Help?}

\author{
\IEEEauthorblockN{Sameel Kazi\IEEEauthorrefmark{1}\textsuperscript{*}, Kapil Joshi\IEEEauthorrefmark{1}, Ali Huzaifa Siddiqui\IEEEauthorrefmark{1}, Samridhi Goel\IEEEauthorrefmark{1}, Yajat Koyande\IEEEauthorrefmark{1}, Adeeb Khan\IEEEauthorrefmark{1}}
\IEEEauthorblockA{\IEEEauthorrefmark{1}Department of Computer Engineering, Sardar Patel Institute of Technology (SPIT), Mumbai, India\\
Email: sameel.kazi25@spit.ac.in}
\thanks{\textsuperscript{*}Corresponding Author. Smart India Hackathon 2026 (SIH26139). Sponsored by Egreen Quanta.}
}

\maketitle

\begin{abstract}
Hybrid quantum-classical machine learning (QML) offers a promising paradigm for computer-aided clinical decision support, but demands rigorous empirical benchmarking against parameter-matched classical baselines. In this study, we design, implement, and systematically benchmark a parameter-compact hybrid Variational Quantum Classifier (VQC) across six clinical diagnostic benchmarks spanning tabular and imaging modalities: Wisconsin Diagnostic Breast Cancer (WDBC, $n=569$), Cleveland Heart Disease ($n=297$), Parkinson's disease dysphonia ($n=195$), Indian Liver Patient Dataset (ILPD, $n=579$), Chronic Kidney Disease ($n=158$), and BreastMNIST ultrasound scans ($n=780$). Benchmarked within a rigorous four-control ablation taxonomy (strictly matched at 73 trainable parameters) across data fractions (10\% to 100\%) and a 10-seed statistical battery with step-down Holm-Bonferroni correction (480 model fits per domain), our hybrid architecture achieves robust statistical parity with classical neural networks across all standard tabular domains ($p \ge 0.105$ post-correction on ILPD; $p \ge 0.078$ overall). Physical hardware executability is confirmed on IBM Quantum's 156-qubit Heron r2 processor (\texttt{ibm\_fez}), achieving 100\% clinical sensitivity (21/21 malignant lesions identified, zero misses) across the locked WDBC test cohort ($n=57$, 1,024 shots/patient, Resilience Level 2 error mitigation, Job ID \texttt{darsoltvr3kc73ejc5i0}) with mean probability divergence $|\Delta p| = 0.0097$ from noiseless simulation. Crucially, the hybrid architecture yields concrete inductive regularization benefits in sample-scarce regimes (Parkinson's, $p=0.00052$, 25-seed replication) and extreme parameter compactness (73 parameters; 292 bytes/round) for communication-efficient federated learning.
\end{abstract}

\begin{IEEEkeywords}
Variational quantum circuits, hybrid quantum-classical machine learning, model ablation, clinical decision support, medical diagnostics.
\end{IEEEkeywords}

\section{Introduction}
Early diagnostic intervention remains the critical determinant of clinical outcomes across oncology, cardiology, and neurology. While parameterized quantum circuits (PQCs) theoretically map biological biomarkers into exponentially high-dimensional Hilbert spaces, recent literature (2024--2026) highlights a reproducibility crisis: published claims of ``quantum advantage'' frequently collapse when evaluated against parameter-matched classical baselines and multiple-testing corrections.

In this work, we present an honest, comprehensive empirical ablation study addressing two fundamental research questions:
\begin{enumerate}
    \item Does variational quantum circuit training outperform an identically parameterized classical neural network?
    \item Under what specific clinical data constraints does training the quantum layer contribute genuine inductive utility over an untrained or unentangled representation?
\end{enumerate}

\section{Methodology and Architecture}

\subsection{Preprocessing and Angle State Preparation}
For a standardized clinical feature vector $\mathbf{x} \in \mathbb{R}^d$, we apply principal component analysis to reduce dimensionality to $q = 6$ orthogonal components. State preparation is executed via Angle Embedding:
\begin{equation}
|\psi(\mathbf{x})\rangle = \bigotimes_{j=0}^{q-1} R_y(x_j) |0\rangle
\end{equation}
requiring a circuit depth of exactly 1 and evading the exponential two-qubit gate overhead of arbitrary amplitude embedding.

\subsection{Variational Quantum Circuit (VQC) Ansatz}
The parameterized quantum state undergoes $L = 2$ variational entangling layers:
\begin{equation}
U(\boldsymbol{\theta}) = \prod_{l=1}^{L} \left( \prod_{j=0}^{q-1} R_z(\phi_{l,j}) R_y(\omega_{l,j}) \right) U_{\text{ent}}
\end{equation}
where $U_{\text{ent}} = \prod_{j=0}^{q-1} \text{CNOT}_{j, (j+1) \bmod q}$ forms a circular entanglement ring. Readout is obtained via Pauli-$Z$ expectation values on all $q$ qubits:
\begin{equation}
m_j = \langle \psi(\mathbf{x}) | U^\dagger(\boldsymbol{\theta}) Z_j U(\boldsymbol{\theta}) | \psi(\mathbf{x}) \rangle \in [-1, 1]
\end{equation}
The final diagnostic probability is computed via a classical linear projection:
\begin{equation}
\hat{y} = \sigma(\mathbf{w}^T \mathbf{m} + b) \in (0, 1)
\end{equation}

\subsection{The Four-Control Experimental Protocol}
To isolate the origin of classification performance, four strictly controlled models are benchmarked:
\begin{itemize}
    \item \textbf{Control A (Hybrid VQC, 73 params):} Full trainable PQC with circular CNOT entanglement.
    \item \textbf{Control B (Classical MLP, 73 params):} Exact parameter-matched classical network ($d \to 8 \to 1$).
    \item \textbf{Control C (Entanglement-Free VQC, 61 params):} Identical to Control A, with all CNOT gates removed.
    \item \textbf{Control D (Untrained Reservoir, 49 params):} Frozen randomized quantum variational rotations.
\end{itemize}

\section{Axiomatic Explainability via Integrated Gradients}
Clinicians require reliable biomarker attributions. We formulate an end-to-end differentiable path integral from quantum expectation values back to raw input biomarkers:
\begin{equation}
\text{IG}_i(\mathbf{x}) = (x_i - x'_i) \times \int_{0}^{1} \frac{\partial F(\mathbf{x}' + \alpha (\mathbf{x} - \mathbf{x}'))}{\partial x_i} \, d\alpha
\end{equation}
By leveraging parameter-shift analytical gradients through the quantum variational layer into the classical preprocessor, our implementation satisfies the \textit{Completeness Axiom}:
\begin{equation}
\sum_{i=1}^{d} \text{IG}_i(\mathbf{x}) = F(\mathbf{x}) - F(\mathbf{x}')
\end{equation}
with empirical relative error strictly bounded below $0.25\%$.

\section{Empirical Results and Discussion}

\subsection{Tabular Parity and Entanglement Ablation}
Across Wisconsin Diagnostic Breast Cancer ($n=569$) and Cleveland Heart Disease ($n=297$), Control A operates at statistical parity with Control B (paired Wilcoxon $p \ge 0.078$ under step-down Holm-Bonferroni correction). Furthermore, deleting entangling gates (Control C) produces zero statistically significant degradation ($p > 0.35$), proving that multi-qubit entanglement is dispensable for low-dimensional tabular clinical manifolds.

\subsection{Inductive Advantage in Scarce Clinical Regimes}
On the low-sample Parkinson's acoustic benchmark ($n=195$), variational quantum training significantly outperforms the untrained reservoir feature map:
\begin{equation}
\Delta \text{AUC} = +2.17\% \pm 0.81\% \quad (p = 0.00052)
\end{equation}
confirmed across 25 pre-registered random seeds, demonstrating that parameterized quantum circuits act as effective regularizers in high-variance, small-$n$ biomedical settings.

\section{Physical IBM Quantum QPU Validation}
We validated inference on IBM Quantum's 156-qubit superconducting processor (\texttt{ibm\_fez}, Heron r2 architecture) using Qiskit 2.3 Runtime with Dynamical Decoupling ($XY4$) and Resilience Level 2 zero-noise extrapolation (ZNE with TREX readout mitigation). Over 1,024 shots per patient across the full locked WDBC test cohort ($n=57$, 21 malignant, 36 benign; Job ID \texttt{darsoltvr3kc73ejc5i0}), the physical hardware achieved 100\% clinical sensitivity (21/21 malignant lesions caught, zero misses) and 100\% specificity at default threshold $\tau=0.50$, matching noiseless statevector simulations with tight fidelity (mean $|\Delta p| = 0.0097$, 100\% prediction agreement). In an exploratory 10-patient cross-domain cohort, uncalibrated hardware sensitivity was 80\% (4/5 caught at $\tau=0.50$), achieving 100\% zero-miss sensitivity when calibrated to clinical triage operating threshold $\tau=0.10$.

\section{Conclusion}
This investigation establishes that while near-term hybrid quantum classifiers achieve parity rather than superiority over classical baselines on tabular clinical data, they provide genuine inductive regularization in sample-scarce regimes, extreme parameter efficiency (73 parameters) for privacy-preserving federated clinical networks, and mathematically sound axiomatic explainability.

\bibliographystyle{IEEEtran}
\begin{thebibliography}{1}
\bibitem{cerezo2021} M. Cerezo et al., ``Variational quantum algorithms,'' \emph{Nature Reviews Physics}, vol. 3, no. 9, pp. 625--644, 2021.
\bibitem{bowles2024} J. Bowles et al., ``Better classical baselines collapse quantum advantage claims,'' \emph{Nature Computational Science}, 2024.
\bibitem{sundararajan2017} M. Sundararajan, A. Taly, and Q. Yan, ``Axiomatic attribution for deep networks,'' in \emph{Proc. ICML}, pp. 3319--3328, 2017.
\end{thebibliography}

\end{document}
